% Folder 79, batch 12 (archivists' pages 221-240).
% Pass: Opus 5.5 (claude-opus-5-5), 2026-10-03 — first pass, unchecked against the pages by a human.
%
\documentclass[11pt,a4paper]{article}
\input{../preamble/grothendieck}

\folder{79}
\batch{12}
\pages{221}{240}
\foldertitle{2-polyèdres réguliers triangulés, et modules libres de rangs 2 sur les anneaux $\Lambda=\mathbb{Z}/n\mathbb{Z}$ et $\Lambda=\mathbb{Z}$ : notes manuscrites (s.d.).}
\dating{[à partir de 1980]}
\watermark{Édition de démonstration}

\begin{document}

\page{221}
\note{Feuillet sans pagination d'auteur, barré en entier de deux longs traits obliques : c'est une première rédaction, abandonnée, de ce que reprennent plus bas les p. 120--122 de l'auteur (pages 236 et 238 du fonds). La phrase initiale continue un texte qui précède le lot.}
Les « conditions d'isomorphie » (la première condition d'homogénéité) dit que l'homomorphisme $G \xrightarrow{\delta} Z$ est surjectif, i.e. qu'on a une suite exacte

(9)\note{le numéro est écrit sur un « 8 ».}

\begin{tikzcd}
1 \arrow[r] & G_0 \arrow[r] & G \arrow[r, "\delta"] & Z \arrow[r] & 1 . \\
 & & Z \arrow[u, hook] \arrow[ur, "\lambda \mapsto \lambda^{2}"'] & &
\end{tikzcd}

\struck{J'aimerais maintenant, dans le contexte}

Je suppose maintenant que la structure de graphe \struck{\ill{}} $(S, A_0)$, correspondant à un $i_0 \in \pi$, est triangulaire et même associée à une structure de 2-polyèdre \add{connexe}, et que $G_0$ soit transitif sur l'ensemble des \uncertain{repères}. Je vais faire bon poids d'hypothèses, quitte à vérifier par la suite, en supposant le 2-polyèdre $P_0$ orientable, soit $G_0^{+}$ le sous-groupe d'indice 2 de $G_0$ qui conserve l'orientation --- on voit que c'est le même pour $P_0$ et pour les $P_i$ ($i \in \pi$). A vrai dire, ce qu'on va utiliser, c'est la donnée d'un sous-groupe
\[
(7) \qquad G_0^{+} \subset G_0
\]
qui ait les propriétés
\[
(8) \qquad
\begin{cases}
\text{a) } G_0^{+} \text{ d'indice 2 dans } G_0, \text{ et invariant dans } G \\
\text{b) } G_0^{+} \text{ est \emph{simplement transitif} sur l'ens } \overrightarrow{A_0} \\
\qquad \text{des arêtes orientées de } P_0 .
\end{cases}
\]

\page{222}
\note{p. 114 de l'auteur.}

\section*{XI) Récapitulation}
\note{Sous le titre, un premier « 4. Récapitulation » est biffé, et le numéro d'abord écrit devant « XI) » est surchargé.}

1. Soit
\[
(1) \qquad \Gamma = (S, S^{\vee}, R \subset S \times S^{\vee})
\]
\note{au-dessus de $S$, $S^{\vee}$, $R$ : « sommets », « cosommets », « incidence ».}
un « prégraphe ».
\marginal{Ex. : \uncertain{bigraphes}, polygraphes, \ill{}, multigraphes, multipolygraphes… (2)}
Deux applications
\[
(2) \qquad \varphi : S \to S^{\vee}, \quad \psi : S \to S^{\vee}
\]
\note{la source de $\psi$ porte un exposant surchargé ; la condition (3), où $t \in S$, demande $\psi : S \to S^{\vee}$.}
sont dites « associées » (ou « $R$-associées ») si on a
\[
(3) \qquad \forall s, t \in S \quad (s, \psi(t)) \in R \iff \text{\struck{$\varphi(s), \psi$}} \ (t, \varphi(s)) \in R .
\]
\struck{\ill{}}

C'est une relation symétrique en $\varphi, \psi$. Si $R$ est séparé en $S^{\vee}$ [i.e.
\[
\check{S} \hookrightarrow \mathfrak{P}(S), \quad s^{\vee} \mapsto {}^{t}R(s^{\vee}) \quad \text{injectif]}
\]
\note{l'exposant de gauche, ici et plus bas (${}^{t}R$, ${}^{t}\varphi$), se lit aussi bien « s » que « t » ; on le transcrit ${}^{t}$, la transposée.}
alors pour tout $\varphi : S \to S^{\vee}$, il existe au plus une $\psi : S \to S^{\vee}$ associée --- \struck{on dit alors que $\varphi$} \add{\ill{}}.

Appelons \emph{quasi-jauges} de $\Gamma$ les $\varphi : S \to S^{\vee}$ telles qu'il existe une $\psi$ associée,
\struck{Dans un sens l'une des quasi-\ill{}}
\struck{\ill{} jauges, ou dans, dans ce cas $\Gamma$ sép. \ill{} quasi-jauge \ill{} une structure de graphe, \ill{} les « arêtes » sont les \ill{} $\{\varphi, \psi\}$, où $\varphi$ et $\psi$ associées. Dans le cas $\Gamma$ séparé en $S^{\vee}$, \ill{}}
\note{passage encadré et biffé, en partie illisible.}
\marginal{en $S^{\vee}$, une \ill{} involution (4) $\varphi \mapsto {}^{t}\varphi$.}

On appelle \emph{jauge} une $\varphi$ associée à elle-même, i.e. telle que
\[
(4) \qquad \forall (s,t) \in S \times S, \quad (s, \varphi(t)) \in R \implies (t, \varphi(s)) \in R
\]
--- ou encore, dans le cas où $\Gamma$ est séparé en $S^{\vee}$
\[
(5) \qquad \varphi = {}^{t}\varphi .
\]
Supposons que \struck{l'on ait} $\varphi$ soit une bijection (on dit alors que $\varphi$ est une \emph{jauge stricte})
\[
(6) \qquad \varphi : S \xrightarrow{\sim} S^{\vee} \quad \text{d'où} \quad \varphi' = \varphi^{-1} : \check{S} \xrightarrow{\sim} S
\]
\note{sous le couple, l'indication « $(\mathrm{Ens}) \times (\mathrm{Ens})$ ».}
et considérons \struck{l'application} \add{l'isomorphisme} déf
\[
(7) \qquad (S, S^{\vee}) \xrightarrow{(\varphi,\ \varphi' \overset{\mathrm{déf}}{=} \varphi^{-1})} (S' \overset{\mathrm{déf}}{=} S^{\vee},\ S'^{\vee} \overset{\mathrm{déf}}{=} S)
\]
on voit que \struck{ce couple de bijections transforme}

\page{223}
\[
(8) \qquad \varphi \times \varphi' \text{ transforme } R \text{ en } R' \overset{\mathrm{déf}}{=} {}^{t}R \quad \text{ssi } \varphi \text{ est une jauge.}
\]
\marginal{Cela revient à dire que si $\varphi : S \rightleftarrows S^{\vee} : \varphi'$, $\varphi' \circ \varphi = \mathrm{id}_S$, alors $(\varphi, \varphi')$ \ill{} $R$ en ${}^{t}R'$ \ill{} ssi $\varphi$ est une jauge.}
Ce qui donne une correspondance biunivoque
\[
(9) \qquad \mathrm{Jauges}_{\text{strictes}}(\Gamma) \simeq
\begin{array}{l}
\text{Ens des isoms } (\varphi, \varphi') : \Gamma \to \Gamma' \\
\text{tels que } \varphi, \varphi' \text{ soient inverses} \\
\text{l'un de l'autre, où } \Gamma' \text{ est le} \\
\text{prégraphe } (S^{\vee}, S, {}^{t}R) \text{ dual de } \Gamma .
\end{array}
\]
Si $\varphi$ est une jauge de $\Gamma$, alors
\[
(10) \qquad \Gamma_{\varphi} = (S, A_{\varphi}), \ \text{où}\ A_{\varphi} \subset \mathfrak{P}(S), \quad A_{\varphi} = \{\{s,t\} \mid (s, \varphi(t)) \in R\}
\]
est un graphe (avec des arêtes qui peuvent être \uncertain{pendantes}), la relation symétrique associée à cette str. de graphe étant
\[
(11) \qquad R_{\varphi} = (\mathrm{id}_S \times \varphi)^{-1}(R) = (\varphi \times \mathrm{id}_S)^{-1}(R) = \{(s,t) \in S \times S \mid \{s, \varphi(t)\} \in R\}
\]
--- c'est le \emph{graphe associé à} $\varphi$. \struck{On trouve ainsi : Si $R$ est}

On trouve ainsi une application
\[
(12) \qquad \mathrm{Jauges}(\Gamma) \longrightarrow
\begin{array}{l}
\text{ens des structures de graphes sur } S \\
\text{telles que l'ens des étoiles } \Sigma_{\gamma} \\
\text{correspondantes soit \struck{égal} \add{contenu dans}} \\
\Sigma = \mathrm{Im}(S^{\vee} \to \mathfrak{P}(S)) \subset \mathfrak{P}(S)
\end{array}
\]
\marginal{$\Sigma \subset \mathfrak{P}(S)$ s'appelle l'ens des \emph{étoiles} du prégraphe, on définit de même $\Sigma^{\vee} \subset \mathfrak{P}(S^{\vee})$, l'ens des coétoiles (\ill{})}
cette application étant \emph{bijective} si $\Gamma$ est séparé en $S^{\vee}$, i.e. si on peut identifier $S^{\vee}$ à un sous-ens de $\mathfrak{P}(S)$
\[
(13) \qquad \check{S} \hookrightarrow \mathfrak{P}(S) .
\]
On trouve \struck{Dans ce cas,} une application \ill{}
\[
(14) \qquad \mathrm{Jauges\ surj}(\Gamma) \longrightarrow
\begin{array}{l}
\text{Ens des str. de graphes } \gamma \text{ sur } S \\
\text{telles que } \Sigma_{\gamma} = \Sigma \\
\text{[i.e. ayant un ens d'étoiles donné]} .
\end{array}
\]
On trouve aussi

\page{224}
\note{p. 115 de l'auteur.}
\[
(15) \qquad \text{Équivalence de groupoïdes}
\]
\[
\begin{array}{l}
\text{groupoïde des graphes} \\
(S, A_0 \subset \mathfrak{P}_{\{1,2\}}(S))
\end{array}
\xrightarrow{\ \approx\ }
\begin{array}{l}
\text{groupoïde des prégraphes} \\
\Gamma = (S, \check{S}, R) \text{ séparés en } S^{\vee} \\
\text{munis d'une jauge \add{surjective} } \varphi
\end{array}
\]
induisant à son tour
\[
(16) \qquad \text{Équivalence de groupoïdes}
\]
\[
\begin{array}{l}
\text{groupoïde des graphes} \\
(S, A_0) \text{ séparés (tout sommet} \\
\text{déterminé par son étoile)}
\end{array}
\xrightarrow{\ \approx\ }
\begin{array}{l}
\text{groupoïde des prégraphes } \Gamma \\
\text{biséparés, munis d'une jauge stricte} \\
\text{(i.e. bijective) } \varphi \text{ (ou encore (9)) d'un} \\
\text{isom } (\varphi, \varphi') : \Gamma \xrightarrow{\sim} \Gamma' \text{ tel que } \varphi, \varphi' \\
\text{inverses l'un de l'autre, avec } \Gamma' \\
\text{prégraphe dual)} .
\end{array}
\]
Soit maintenant
\[
(17) \qquad \mathcal{A} = \mathrm{Aut}(\Gamma) = \Bigl\{\lambda = (\lambda_S, \lambda_{\check{S}}) \in \mathfrak{S}_S \times \mathfrak{S}_{\check{S}} \ \Bigm|\ 
\begin{array}{l}
\forall (s, \check{s}) \in S \times \check{S} \\
(s, s^{\vee}) \in R \iff (\lambda_S s, \lambda_{S^{\vee}} s^{\vee}) \in R
\end{array}
\Bigr\}
\]
\struck{et soit $\varphi$ une jauge, d'où une structure de graphe $\gamma_{\varphi} = (S, A_{\varphi})$. Alors (18) conditions équivalentes : a) $\lambda_{\check{S}} \varphi = \varphi \lambda_S^{-1}$ ; b) $\lambda_S$ est \ill{}}
\note{passage encadré et biffé de traits obliques.}
Notons que
\[
(18) \qquad
\begin{array}{l}
\text{Si } \Gamma \text{ est} \\
\text{séparé en } S^{\vee} \\
\text{i.e. } S^{\vee} \hookrightarrow \mathfrak{P}(S) \\
\text{alors}
\end{array}
\quad
\begin{array}{l}
\lambda \mapsto \lambda_S \\
\mathcal{A} \to \mathfrak{S}_S
\end{array}
\quad
\begin{array}{l}
\text{induit un \struck{bijection} \add{isom.} du groupe} \\
\mathcal{A} \text{ avec le groupe des} \\
\text{permutations } f \text{ de } S \\
\text{telles que } f(\Sigma) = \Sigma \text{ (tr. de str.) i.e.} \\
\text{qui transforment étoile en} \\
\text{étoile,}
\end{array}
\]
où $\Sigma = \mathrm{Im}(S^{\vee} \to \mathfrak{P}(S))$.

Soient maintenant $\varphi$ une jauge surjective de $\Gamma$ (séparé en $S^{\vee}$), \struck{définissant donc} \add{correspondant à} une structure de graphes
\[
(19) \qquad \gamma = \gamma_{\varphi} = (S, A_{\varphi}) \quad \text{avec} \quad \underbrace{\Sigma_{\gamma}}_{\text{ens des étoiles de } \gamma} = \Sigma
\]
\struck{On a alors \ill{}} et soit $f \in \mathfrak{S}_S$. On a alors

\page{225}
\[
(20) \qquad \text{On a équivalence}
\quad
\begin{array}{l}
f \text{ est une \emph{homothétie} pour } \gamma_{\varphi} \\
\text{i.e. } \forall (s,t) \in S \times S, \\
\qquad \{s,t\} \in A_{\varphi} \iff \{f^{-1}(s), f(t)\} \in A_{\varphi} \\
\text{i.e. } \forall t \in S, \\
\qquad f(\mathrm{Et}_{A_{\varphi}}(t)) = \mathrm{Et}_{A_{\varphi}}(f^{-1}(s))
\end{array}
\]
\[
\iff
\begin{array}{l}
\exists \lambda \in \mathcal{A}, \text{ avec } f = \lambda_S, \\
\text{et} \\
\lambda_{\check{S}} \varphi = \varphi \lambda_S^{-1}
\end{array}
\]
\note{la dernière ligne de gauche porte bien « $\forall t$ » puis « $f^{-1}(s)$ ».}
On trouve ainsi
\[
(21) \qquad
\begin{array}{l}
\text{Correspondance} \\
\text{biunivoque} \\
\text{entre}
\end{array}
\quad
\begin{array}{l}
\text{homothéties} \\
\text{du graphe} \\
\gamma_{\varphi} = (S, A_{\varphi})
\end{array}
\ \simeq\ 
\begin{array}{l}
\text{Automorphismes } \lambda \in \mathcal{A} \\
\text{du prégraphe } \Gamma, \text{ tels que} \\
\lambda_{\check{S}} \varphi = \varphi \lambda_S^{-1} .
\end{array}
\]
On dit qu'une jauge $\varphi$ \struck{\ill{}} (on suppose ici $\Gamma$ séparé en $S^{\vee}$ et $\varphi$ surjective) est compatible avec $\lambda$, ou est une \emph{$\lambda$-jauge}, si elle satisfait
\[
(22) \qquad \lambda_{\check{S}} \varphi = \varphi \lambda_S^{-1} \quad \text{i.e.} \quad \lambda_{\check{S}} \varphi \lambda_S = \varphi \quad \text{\struck{\ill{}}}
\]
(dans ce cas, si $Z \subset \mathcal{A}$, on définit les $Z$-jauges comme les jauges satisfaisant (22) $\forall \lambda \in Z$).
\marginal{NB Si $\varphi$ bijective, la condition (22) équivaut à (22') $\varphi(\lambda) = \lambda_{\check{S}}^{-1}$, (22'') $\varphi^{-1}(\lambda_{\check{S}}) = \lambda_S^{-1}$.}
\note{lecture incertaine des deux formules marginales, telles qu'écrites.}

Ainsi, si $\varphi$ est une $\lambda$-jauge, alors $\lambda_S$ est une homothétie pour $\gamma_{\varphi}$, et inversement dans le cas où $\Gamma$ est séparé en $\check{S}$.

\struck{Supposons} On a
\[
(23) \qquad
\begin{array}{l}
\text{Si } \varphi \text{ est une \add{bijective} } \{\lambda,\mu\}\text{-jauge, où } \lambda, \mu \in \mathcal{A}, \\
\text{alors } \varphi \text{ est une } \lambda\mu\text{-jauge ssi } \lambda\mu = \mu\lambda
\end{array}
\]
ce qui implique que \struck{pour} \add{deux homothéties d'}un graphe séparé, \struck{\ill{}} \add{bijectives}, \struck{\ill{}} homothéties ssi elles commutent. Si $Z$ est un sous-groupe de $\mathcal{A}$, tel qu'il existe une $Z$-jauge bijective, alors $Z$ doit être commutatif.

Donnons-nous maintenant \struck{\ill{}}
\[
(24) \qquad Z \subset \mathcal{A} = \mathrm{Aut}(\Gamma) \quad \text{un sous-groupe \emph{commutatif}}
\]
et considérons \struck{\ill{} $Z$-jauge $\varphi$.}

\page{226}
\note{p. 116 de l'auteur.}
\[
(25) \qquad \sqcap_{\hat{Z}} = \text{ens des $Z$-jauges \struck{bijectives} \add{strictes} de } \Gamma
\]
alors $Z$ opère sur $\sqcap_{\hat{Z}}$ par
\[
(26) \qquad T_{\lambda} \varphi \overset{\mathrm{def}}{=} \lambda_{\check{S}} \circ \varphi = \varphi \circ \lambda_S^{-1}
\]
\note{les deux formules sont numérotées « (25) » ; la seconde est ici numérotée (26), comme l'appelle la suite (27). Le symbole $\sqcap$ est rendu tel qu'il est dessiné ; son indice se lit $\hat{Z}$ ici, $Z$ plus bas.}
--- l'opération correspondante sur l'ens des graphes associés aux $\varphi$ étant donnée par
\[
(27) \qquad T_{\lambda} A_{\varphi} = \{\{s, \lambda_S(t)\} \in \mathfrak{P}(S) \mid s, t \text{ t.q. } \{s,t\} \in A_{\varphi}\} .
\]
Si $\varphi$ est surjective, alors pour $\varphi$ fixé, la connaissance de $T_{\lambda}\varphi = \lambda_{\check{S}} \circ \varphi$ implique celle de $\lambda_{\check{S}}$, si $\varphi$ est injective elle implique celle de $\lambda_S$, donc si $\varphi$ est bijection, la connaissance de $T_{\lambda}\varphi$ implique celle de $\lambda$, i.e. \struck{l'opération}
\[
(27) \qquad
\begin{array}{l}
\text{sur l'ens } \sqcap_Z \text{ des $Z$-jauges \emph{bijectives}, $Z$ opère} \\
\text{\emph{librement} --- donc sur \struck{l'ens des graphes \ill{}} l'ens des} \\
\text{structures de graphes \emph{séparés} sur un ens } S, \\
\text{\emph{commutant} avec un sous-groupe commutatif } Z \subset \mathfrak{S}_S \\
\text{comme groupe d'homothéties, } Z \text{ opère} \\
\text{\emph{librement}.}
\end{array}
\]
\note{ce second « (27) » est du manuscrit.}
Notons
\[
(28) \qquad \mathcal{G} = \mathrm{Cent}_{\mathcal{A}}(Z) \text{ \struck{\ill{}}}, \quad \text{d'où} \quad Z \subset \mathcal{G} \subset \mathcal{A} .
\]
On trouve
\[
(29) \qquad
\left\{
\begin{array}{l}
\text{Soit } \varphi_0 \in \sqcap_Z \text{ une $Z$-jauge stricte (i.e. bij.) de } \Gamma, \ \text{\struck{alors}} \text{ soit} \\
\text{\struck{$\mathcal{G}_{\varphi_0}$}} \ \mathcal{G}_{\varphi_0} = \{g \in \mathcal{G} \mid g_{\check{S}} \varphi = \varphi\, g_S^{-1} \text{ i.e. } \varphi_0 \text{ est une $g$-jauge}\} \\
\text{(NB } \mathcal{G}_{\varphi_0} \text{ n'est pas nécessairement un sous-groupe de } \mathcal{G}) \\
\text{alors} \\
\qquad g \longmapsto T_g \varphi \overset{\mathrm{def}}{=} g_{\check{S}} \varphi_0 = \varphi_0 g_S^{-1} \\
\text{est une bijection } \mathcal{G}_{\varphi_0} \xrightarrow{\ \sim\ } \sqcap_Z
\end{array}
\right.
\]
\note{dans la marge gauche de (29), un mot biffé illisible ; l'indice « $Z\varphi$ » sous « bijection » est surchargé.}
Mais on fera attention que par (23), $\mathcal{G}_{\varphi}$ est un ss-groupe de $\mathcal{G}$ ssi il est formé d'éléments 2 : 2 permutables --- auquel cas c'est un sous-groupe commutatif, \struck{Dans ce}
\note{la phrase continue p. 117 de l'auteur.}

\page{227}
\note{feuillet sans numéro d'auteur, entre ses p. 116 et 117.}
cas et dans ce cas seulement, on peut considérer que $\hat{Z} = \mathcal{G}_{\varphi_0}$ opère sur $\sqcap_Z$ (et ce seulement sur $\varphi_0$), --- et alors $\sqcap_Z = \sqcap_{\hat{Z}}$ est un torseur sous $\hat{Z}$.

Nous nous intéressons, pour $(\Gamma, Z \subset \mathcal{A}_{\Gamma} = \mathrm{Aut}(\Gamma))$ donné, avec $\Gamma = (S, S^{\vee}, R)$, $Z$ ss-groupe commutatif, \struck{aux} à la donnée d'un sous-torseur $\pi$ de $\sqcap_Z$, sous l'opération de $Z$. Dans le cas d'une $\Gamma$ biséparée, cela signifie qu'on s'intéresse aux « multigraphes » séparés à groupe d'homothéties $Z$ :
\[
(30) \qquad \underbrace{\pi}_{\text{sous $Z$-torseur}} \subset \sqcap_Z = \text{ens des $Z$-jauges bijectives} .
\]
On trouve alors des bijections canoniques
\[
(31) \qquad S^{\vee} \simeq \pi \wedge^{Z} (S!) , \qquad S \simeq \check{\pi} \wedge^{Z} (\check{S}!)
\]
\note{lecture de « $\pi \wedge^{Z}$ » et de « $(S!)$ » incertaine.}
dont la deuxième peut être considérée, soit comme une \uncertain{variante} de la première, soit comme un \uncertain{cas} de la première, \uncertain{par} les isomorphismes \uncertain{de} \uncertain{dualité} : la situation duale
\[
(32) \qquad (\underset{\substack{\| \\ (S^{\vee}, S, {}^{t}R)}}{\Gamma'},\ \underset{\substack{\| \\ {}^{t}Z = \{\{\lambda_{\check{S}}, \lambda_S\} \mid \lambda \in Z\}}}{Z'},\ \pi' \text{ \struck{\ill{}}}) \qquad \pi' = \{\varphi^{-1} \mid \varphi \in \pi\}
\]
on voit que
\[
\varphi \longmapsto \varphi^{-1} \quad \pi \longrightarrow \pi'
\]
\struck{\ill{}} satisfait à $(T_{\lambda}\varphi)^{-1} = T'_{\lambda^{-1}}(\varphi^{-1})$ et définit donc un isom canon. de $Z$-torseurs
\[
(33) \qquad \pi' \simeq \pi^{[-1]} \overset{\mathrm{def}}{=}
\begin{array}{l}
\pi \text{ avec opér. de } Z \text{ opposée de} \\
\text{l'opération donnée } T_{\lambda}, \text{ i.e.} \\
\lambda \text{ opérant par } T_{\lambda^{-1}} \ (\overset{\mathrm{déf}}{=} T'_{\lambda}) .
\end{array}
\]

\page{228}
\note{p. 117 de l'auteur. Le symbole noté ici $\flat$ en exposant ($S^{\flat}$, $R^{\flat}$) est une sorte de bémol ; lecture incertaine du signe.}
\struck{L'opération de $Z$}
Considérons
\[
(34) \qquad S^{\flat} = S / Z_S , \quad S^{\vee\flat} = S^{\vee} / Z_{S^{\vee}}
\]
alors les $\varphi \in \pi$ définissent par passage aux quotients une même application bijective
\[
(35) \qquad \varphi^{\flat} : S^{\flat} \xrightarrow{\ \sim\ } S^{\vee\flat}
\]
par quoi on identifie $S^{\flat}$, $S^{\vee\flat}$ : un même ensemble. \struck{\ill{}} Les orbites \struck{\ill{}} \add{$Z$-orbites} de $Z_S$ dans $S$ s'appellent \emph{\uncertain{rayons}}, \struck{\ill{}} celles dans $S^{\vee}$ s'appellent \uncertain{corayons}, --- donc la donnée de $\pi$ nous définit une bijection entre \add{ens. des} rayons et \add{ens. des} corayons. On pose
\[
(36) \qquad R^{\flat} = R / Z \quad \text{\struck{et l'application}}
\]
et l'injection $R \hookrightarrow S \times S^{\vee}$ définit par passage aux quotients
\[
(37) \qquad R^{\flat} \longrightarrow S^{\flat} \times S^{\vee\flat}
\]
qui n'est plus injective en général, \struck{en fait,} mais qu'il l'est \struck{\ill{} sur la \ill{}} \struck{\ill{} $Z$ sur $R$ provenant d'une opération de $Z \times Z$ sur $S \times S^{\vee}$} \add{\ill{}} souvent. Supposons pour simplifier que $Z$ opère librement sur $S$ et $S^{\vee}$ (\ill{} l'une d'ailleurs \ill{} implique \ill{} $\pi$…). Alors les rayons et les corayons sont des torseurs sous $Z$.
\marginal{hypothèse vérifiée dans le cas bisép. « homogène »}
\note{sous la note marginale, un bloc biffé de hachures.}
Soient \struck{\ill{}}
\[
D \subset S , \quad \check{D} \subset S^{\vee} , \qquad D \in S^{\flat} , \ \check{D} \in S^{\vee\flat}
\]
\note{les deux appartenances sont écrites sous les inclusions, avec le signe $\in$ couché.}
\struck{\ill{}} rayons et corayons « liés » i.e. \struck{\ill{}} correspondant : un $(s^{\flat}, s^{\vee\flat}) \in S^{\flat} \times S^{\vee\flat}$ dans l'image de $R^{\flat}$, i.e. tels que
\[
(38) \qquad \exists s_0 \in D, \ s_0^{\vee} \in \check{D} \text{ tels que } (s_0, s_0^{\vee}) \in R
\]
(NB $D = Z s_0$, $\check{D} = Z s_0^{\vee}$). Alors la \struck{partie de} \add{fibre de} $R^{\flat}$ \struck{\ill{}} $(s^{\flat}, s^{\vee\flat})$ est le quotient

\page{229}
par $Z$ de la partie de $R$, stable sous $Z$, \struck{définie} \add{définie par}
\[
R(s^{\flat}, s^{\vee\flat}) = \Bigl\{(s, s^{\vee}) \in R \ \Bigm|\ 
\begin{array}{l}
s \text{ au-dessus de } s^{\flat} \\
s^{\vee} \text{ au-dessus de } s^{\vee\flat}
\end{array}
\Bigr\}
\]
\note{devant $R(s^{\flat}, s^{\vee\flat})$, un « $Z$ » biffé.}
\[
(39) \qquad
\begin{array}{c}
R(s^{\flat}, s^{\vee\flat}) \\
\uparrow \ (\lambda,\mu) \mapsto (\lambda s_0, \mu s_0^{\vee}) \\
E(s_0, s_0^{\vee}) \overset{\mathrm{def}}{=} \Bigl\{(\lambda,\mu) \in Z \times Z \ \Bigm|\ 
\begin{array}{l}
(\lambda s_0, \mu \check{s}_0) \in R \text{ i.e.} \\
((\mu^{-1}\lambda) s_0, \check{s}_0) \in R \text{ i.e.} \\
\mu^{-1}\lambda s_0 \in \underbrace{E(\check{s}_0)}_{\text{étoile dans $S$ définie par } \check{s}_0}
\end{array}
\Bigr\} .
\end{array}
\]
\note{la flèche verticale porte à sa gauche un « 2 » ou un « Z ».}
On trouve ainsi une bijection can. $Z(s_0, s_0^{\vee})$,
\[
(40) \qquad \text{fibre de } R^{\flat} \text{ en } (s^{\flat}, s^{\vee\flat}) \simeq
\begin{array}{l}
\text{\add{ens. des}} \ z \in Z \text{ \struck{sous-groupe des}} \\
\text{tels que } (z s_0, \check{s}_0) \in R \\
\text{i.e. } (s_0, z^{-1}\check{s}_0) \in R \\
\text{i.e. } z s_0 \in E
\end{array}
\]
Cette bijection, \add{qui dépend du choix de $s_0, s_0^{\vee}$}, en fait, ne dépend pas de l'hypothèse $(s_0, \check{s}_0) \in R$ --- les deux membres de (40) pouvant être vides, \struck{\ill{}} l'hypothèse $(s_0, s_0^{\vee}) \in R$ signifie simplement que $Z_{s_0, \check{s}_0} \ni 1$. Quand on change de $(s_0, \check{s}_0)$, le sous-ens. en question est remplacé par un translaté. \struck{Dans le cas où il existe une $Z$-jauge stricte $\varphi$, on a l'équivalence}
On a l'équivalence [si $Z$ opère librement sur $S$, $S^{\vee}$]
\[
(41) \qquad
\text{\struck{équivalence}}\ R^{\flat} \to Z^{\flat} \times Z'^{\flat} \text{ injective}
\iff
\begin{array}{l}
\forall E \in \Sigma,\ s_0 \in E, \\
\lambda \in Z,\ \lambda s_0 \in E \implies \lambda = 1
\end{array}
\]
\note{« $Z^{\flat} \times Z'^{\flat}$ » tel qu'écrit, là où l'on attend $S^{\flat} \times S^{\vee\flat}$.}
\[
\underset{\text{\uncertain{quand on se donne} une $Z$-jauge stricte}}{\iff}
\forall (s,t) \in S \times S, \ \bigl(\{s,t\} \in A_0 \text{ et } \{s, \lambda t\} \in A_0 \implies \lambda = 1\bigr)
\]
\[
\iff \forall \lambda \in Z, \ \lambda \neq 1, \ \ldots\ A_0 \cap A_{\lambda} = \emptyset
\]
\note{à gauche, un passage encadré, biffé de traits obliques, reprend les mêmes équivalences ($\forall \lambda \in Z, \lambda \neq 1$, $A_0 \cap A_{\lambda} = \emptyset$…).}
\emph{Cette condition signifie que} si $D$ et $D^{\vee}$ sont \struck{deux} \add{un} rayon et corayon incidents, alors pour tout $s \in D$, $\exists!\, s^{\vee} \in D^{\vee}$ tel que $s, \check{s} \in D^{\vee}$\note{ainsi sur la page ; on attend $(s, \check{s}) \in R$.}. Donc si $s$ remplacé par $\lambda s$, $s^{\vee}$ remplacé par $\lambda s^{\vee}$. Donc \struck{(41) équivaut à :} \struck{\ill{}}
\[
(41') \qquad
\begin{array}{l}
\text{Pour deux rayons } D \in S^{\flat}, \ D^{\vee} \in S^{\vee\flat} \text{ incidents,} \\
\text{la relation \struck{\ill{}}} \ (s, \check{s}) \in R \text{ est} \\
\text{une relation bijective, établissant un isom. de $Z$-torseurs} \\
\qquad D \simeq \check{D}
\end{array}
\]

\page{230}
\note{p. 118 de l'auteur.}
L'interprétation en termes d'une \struck{\ill{}} $Z$-jauge stricte, i.e. d'une structure de graphe $(S, A_0)$ ayant $Z$ comme groupe d'homothéties, est un peu différente :
\[
(41'') \qquad
\begin{array}{l}
\text{(41) signifie que pour deux rayons } D, D' \in S^{\flat} \text{ incidents} \\
\text{(i.e. tels qu'il existe } (s,s') \in D \times D' \text{ avec } (s,s') \in A_0\text{),} \\
\forall s \in D, \ \exists!\, s' \in D' \text{ adjacent à } s \text{ --- cette bijection } D \simeq D' \\
\text{est un \emph{anti}-isomorphisme pour les structures de torseur,} \\
\text{donc établit un iso.} \\
\qquad D \wedge^{Z} D' \simeq \mathbf{1}_Z
\end{array}
\]
Enfin, en termes de la donnée d'un $\pi$, sous $Z$-torseur de $\sqcap_Z$, on a une formulation un peu différente encore. Introduisons
\[
(42) \qquad A = A_{\pi} = \bigcup_{i \in \pi} A_i \subset \mathfrak{P}_{\{1,2\}}(S)
\]
qui définit une structure de graphe sur $S$, on a alors
\[
\text{\struck{(43) Pour deux rayons $D, D'$ de $S$ « $\pi$-incidents » i.e. …}}
\]
\note{un passage encadré est biffé d'une grande diagonale : « Pour deux rayons $D, D'$ de $S$, « $\pi$-incidents » i.e. … $s'^{\flat} = \varphi^{\flat}(s^{\flat})$ (35), ou $s^{\flat} = \varphi^{\flat}(s'^{\flat})$, ssi $\exists s \in D, s' \in D'$ tels que $\{s,s'\} \in A$, on a ssi $\forall s \in D, s' \in D'$, … $\{s,s'\}$ … ».}
Soient \struck{\ill{}} $D, D' \subset S$, \struck{deux} rayons, correspondant à $s^{\flat}, s'^{\flat} \in S^{\flat}$, conditions équivalentes
\[
(43) \qquad
\left|
\begin{array}{l}
\text{a) } s'^{\flat} = \varphi^{\flat}(s^{\flat}) \\
\text{b) } s^{\flat} = \varphi^{\flat}(s'^{\flat}) \\
\text{c) } \forall (s,s') \in D \times D', \ \ldots \ \{s,s'\} \in A \quad (42) \\
\text{d) } \exists (s_0, s'_0) \in D \times D' \ \ldots \ \{s_0, s'_0\} \in A
\end{array}
\right.
\]
\note{en c) et d), quelques mots courts illisibles avant $\{s,s'\}$ et $\{s_0,s'_0\}$.}
On dit alors que $D, D'$ sont \add{$\pi$-}\emph{incidents}, \struck{\ill{}} \add{conditions} qui ne dépend d'ailleurs que de la donnée de $Z_S \subset \mathfrak{S}_S$ (permettant de définir $S^{\flat}$ et la notion de rayon),
\note{la phrase se poursuit p. 119 de l'auteur.}

\page{231}
et de \struck{$A$} structure de graphe $(S, A)$, et non de la famille $\pi$ de sous-ens $A_i \subset A$, $i \in \pi$. Cela posé, on a
\note{un passage encadré, numéroté (44) et (45), est barré de deux grandes diagonales ; on le transcrit tel qu'il se lit.}
\struck{(44) Les cond. équivalentes (41) (41') (41'') signifient aussi que pour deux $Z$-rayons incidents $D, D' \subset S$, et pour $\forall (s,s') \in D \times D'$, il existe un unique $i = i_{s,s'} \in \pi$ tel que $\{s,s'\} \in A_i$, soit \ill{}. On a $i_{\lambda s, s'} = i_{s, \lambda s'}$, et on trouve ainsi un iso. canonique}
\[
\text{\struck{$(45) \qquad D \wedge^{Z} D' \simeq \pi$}}
\]
\struck{caractérisé par $i_{s,s'} = \tau(s \wedge^{Z} s')$, i.e. par la relation \ill{} $s,s'$.}

(44) Les conditions équivalentes (41), (41') (41'') signifient aussi que les $A_i$ sont mutuellement disjoints, \add{donc} \struck{\ill{}} on a une application canonique « type »
\[
A \xrightarrow{\ \tau_A\ } \pi .
\]
Si \struck{alors} $D, D'$ sont deux rayons incidents, \struck{\ill{} (45) Si cette \ill{}}
\[
(s,s') \longmapsto \tau_A(\{s,s'\}) \qquad D \times D' \xrightarrow{\ \tau_{D,D'}\ } \pi
\]
satisfait
\[
(46) \qquad \text{\struck{$\tau_A(\{s,s'\})$}} \ \tau_{D,D'}(\lambda s, \mu s') = \lambda\mu\, \tau_{D,D'}(s,s') \ \text{\struck{$= \tau_{D,D'}(s, \lambda s')$}} ,
\]
et se réduit donc en un isom. de $Z$-torseurs
\[
(47) \qquad D \wedge D' \xrightarrow{\ \sim\ } \pi .
\]
\marginal{donc la structure de multigraphe $(A_i)_{i \in \pi}$ sur $S$ \ill{} se \ill{} \ill{} (S) … est connue, connaît $S$, $A = \coprod_{i \in \pi} A_i \subset \mathfrak{P}_{\{1,2\}}(S)$, et \ill{} $\pi$ \ill{} (\ill{} \ill{} de $Z$…) \ill{} de la condition que $Z \subset \mathfrak{S}_S$ \ill{} groupe d'homothéties de $(S, A_i)$, \ill{} \ill{} (45) \ill{}. Dans le cas \ill{}}
\note{note marginale écrite en oblique dans le coin gauche, très serrée et en partie surchargée ; lecture fragmentaire.}

\emph{NB} La donnée d'un prégraphe $\Gamma$, d'un groupe $Z$ commutatif d'automorphismes --- pouvant être vu comme groupe d'homothéties virtuelles --- \uncertain{provenant de} \ill{} défini dans $S$ (disons, \ill{} plus $S^{\vee}$)
\note{la phrase se poursuit p. 120 de l'auteur ; ce feuillet ne porte pas de numéro d'auteur.}

\page{232}
\note{p. 119 de l'auteur (numéro surchargé sur « 118 »).}
deux types de sous-ensembles remarquables, d'une part les \emph{étoiles} (indépendantes de la donnée de $Z$), d'autre part les \emph{rayons} (en dépendant par contre par le biais de la donnée de $Z_S \subset \mathfrak{S}_S$). La relation \add{d'intersection} des rayons entre eux est évidente (deux rayons distincts ont une intersection vide), alors que celles des deux étoiles peut-être plus complexe (il faudrait l'étudier systématiquement \uncertain{$g$}, dans des cas où les données structurales sont assez riches pour être pas \struck{\ill{}} trop pauvres, p.ex. sur un $\Lambda$-module libre de rang $2$…). La relation entre étoiles et rayons est simple, dans le cas justement de l'hyp. simplificatrice (41) et reformulations (41'), (41''), (44).
\struck{Elle signifie que} \struck{Notons} \struck{ces conditions,} Notons \add{en effet} que si $s \in S$, $\check{s} \in \check{S}$, et si $E \subset S$ l'étoile définie par $\check{s}$, alors
\[
s \in E \iff (s, \check{s}) \in R
\]
donc \struck{l'hyp.}
\[
(48) \qquad
\begin{array}{l}
\text{\add{La condition}} \ \text{\struck{L'hyp.}} \ (41) \text{ signifie aussi que l'intersection dans $S$} \\
\text{d'une étoile $E$ et d'un rayon $D$ est, soit vide,} \\
\text{soit réduite à un point.}
\end{array}
\]
Dans le cas où on part d'une structure de graphe, et où $E$ est l'étoile de $s \in E$, alors $E \cap D \neq \emptyset$ ssi les rayons $D$ et \struck{\ill{}} $D' = Z.s$ \emph{sont incidents}.
\note{les deux derniers mots sont reportés sur la bande qui attache au bas du feuillet la feuille « 119 bis » (page suivante).}


\page{233}
\note{feuille numérotée « 119 bis » (cerclé) par l'auteur, attachée par une bande collée au bas de la p. 119 ; c'est une note introduite par « NB ». La numérotation (48.1)–(48.6) est surchargée, sur « (49.…) » semble-t-il.}
\emph{NB} La notion de prégraphe $\Gamma = (S, S^{\vee}, R)$ avec groupe $Z$ d'homothéties virtuelles
\[
(48.1) \qquad Z \subset \mathrm{Aut}(\Gamma) , \quad Z \text{ ss-gr. comm.}
\]
satisfaisant les conditions
\[
(48.2) \qquad
\left\{
\begin{array}{l}
\text{a) } Z \text{ opère librement sur $S$ et sur } S^{\vee} \\
\text{b) \struck{\ill{}} On a la condition (41) (ou (41'))}
\end{array}
\right.
\]
peut s'expliciter ainsi. On peut identifier, à cause de a), $S$ et $\check{S}$ comme des $Z$-torseurs relatifs sur deux ens $S^{\flat}$ et $\check{S}^{\flat}$, liés par une relation
\[
(48.3) \qquad R^{\flat} \subset \underbrace{S^{\flat}}_{\text{ens des rayons}} \times \underbrace{\check{S}^{\flat}}_{\text{ens des corayons}}
\]
\[
(R^{\flat} \text{ « relation d'incidence » entre rayons et corayons})
\]
\note{sous $R^{\flat}$, les deux projections $p : R^{\flat} \to S^{\flat}$ et $\check{p} : R^{\flat} \to S^{\vee\flat}$ ; « relation d'incidence » récrit au-dessus d'un premier « relation d'incidence » biffé.}
et on a un isom. de $Z$-torseurs
\[
(48.4) \qquad p^{*}(S/S^{\flat}) \simeq \check{p}^{*}(\check{S}/S^{\vee\flat})
\]
--- la relation $R$ entre $S$ et $\check{S}$ s'identifie : $p^{*}(S/S^{\flat})$ et $\check{p}^{*}(\check{S}/S^{\vee\flat})$ identifiés par (48.4), s'envoyant dans $S$ et dans $S^{\vee}$ via les diag. commutatifs

(48.5)

\begin{tikzcd}[column sep=small, row sep=small, nodes={font=\scriptsize}]
S \arrow[d] & p^{*}(S/S^{\flat}) \arrow[l] \arrow[dr] & \simeq & \check{p}^{*}(\check{S}/S) \arrow[dl] \arrow[r, leftrightarrow] & \check{S} \arrow[d] \\
S^{\flat} & & R^{\flat} \arrow[ll, "p_{\flat}"] \arrow[rr, "\check{p}"'] & & S^{\vee\flat}
\end{tikzcd}

\note{diagramme surchargé : sous $p^{*}(S/S^{\flat})$ une flèche verticale vers un symbole raturé (« $R'^{\flat}$ » ?) ; on garde les flèches lisibles. L'argument « $(\check{S}/S)$ » est tel qu'écrit.}

Ceci posé, \struck{\ill{}} la donnée d'une \emph{$Z$-jauge \uncertain{stricte}} \struck{\ill{}} \struck{torseur $\pi$ de $\sqcap_Z$} peut s'interpréter en termes des
\[
(48.6) \qquad
\left\{
\begin{array}{l}
\text{i) une \struck{\ill{}} \add{bijection} } S^{\flat} \xrightarrow{\ \varphi^{\flat}\ } S^{\vee\flat} \text{ qui soit une $R^{\flat}$-jauge} \\
\qquad \text{(i.e. donnant lieu à une relation symétrique sur } S^{\flat} \times S^{\flat}, \\
\qquad \text{soit \struck{\ill{}}} \ R^{\flat}_{S^{\flat}}) \\
\text{ii) un \struck{\ill{}} \add{isomorphisme}} \\
\qquad (S/S^{\flat}) \wedge^{Z_{S^{\flat}}} (\varphi^{\flat *}(S^{\vee}/S^{\vee\flat})) \simeq \mathbf{1}_{Z_{S^{\flat}}}
\end{array}
\right.
\]
\struck{et la donnée d'une}
\note{dans la marge gauche, à hauteur de (48.6), une variante biffée de traits obliques (« d'un \ill{} : 2 données \ill{} … $Z$, … $\lambda^{-1} \circ \varphi$ ») ; on ne la transcrit pas plus avant.}

\page{234}
\note{suite de la feuille « 119 bis » (numéro répété dans la marge gauche) ; au-dessus de la bande collée, quelques mots coupés du bas de la feuille voisine, non transcrits.}
\struck{\ill{}} graphe, avec groupe d'homothéties $Z$ opérant librement, et satisfaisant la condition (41''), \struck{équivaut à celle} revient à la donnée \struck{des choses} suivantes
\[
(48.7) \qquad
\left\{
\begin{array}{l}
\text{a) une \struck{\ill{}} \ill{} $S^{\flat}$, une relation symétrique } R^{\flat} \text{ (entre } p_1, p_2 : R^{\flat} \rightrightarrows S^{\flat}) \\
\qquad \text{(i.e. un \emph{graphe} } \gamma^{\flat}) \\
\text{b) un $Z$-torseur relatif } S/S^{\flat} \text{ sur } S^{\flat} \\
\text{c) un isomorphisme} \\
\qquad p_1^{*}(S/S^{\flat}) \wedge p_2^{*}(S/S^{\flat}) \simeq \underbrace{\mathbf{1}_{Z_{S^{\flat}}}}_{\text{torseur trivial}} .
\end{array}
\right.
\]
\marginal{c'est \struck{\ill{}} la relation symétrique d'une structure de graphe \ill{}}
\note{note marginale en oblique, en partie biffée.}

Enfin, dans cette optique, la donnée d'un multigraphe $\pi$, \struck{\ill{}} \struck{$\pi$} avec groupe d'homothéties, revient à donner \ill{} les données suivantes
\[
(48.8) \qquad
\left\{
\begin{array}{l}
\text{a) Un graphe } \gamma^{\flat} = (S^{\flat}, \underbrace{R^{\flat} \subset S^{\flat} \times S^{\flat}}_{\text{rel. symm.}}) \\
\text{b) Un $Z$-torseur relatif } S/S^{\flat} \\
\text{c) Un $Z$-torseur absolu } \pi \\
\text{d) un isomorphisme} \\
\qquad p_1^{*}(S/S^{\flat}) \wedge p_2^{*}(S/S^{\flat}) \simeq \pi
\end{array}
\right.
\]
Ceci : la description (48.7), \add{ou} $i \in \pi$ définit sous une structure de graphe à homothéties sur $S$, \add{les sommets}, puisqu'il définit un iso $p_1^{*}(S/S^{\flat}) \wedge p_2^{*}(S/S^{\flat}) \simeq \pi \xrightarrow[\text{via } i]{\ \sim\ } \mathbf{1}$.
\note{« (44.7) » sur la page, là où l'on attend (48.7) ; la phrase est elliptique et le début se lit mal.}

\emph{Ex.} Si $M$ module libre de rg $2$ sur $\Lambda$, $S = M^{*}/\pm 1$, $\check{S}^{*} \simeq \check{M}^{*}/\pm 1$, $Z = \Lambda^{*}/\pm 1$, \struck{\ill{}} $\pi = (\det M)^{*}/\pm 1$, on a
\[
S^{\flat} = M^{*}/\Lambda^{*} \subset \mathbb{P}(M)(\Lambda) \quad \text{(égalité si } \mathrm{Pic}(\Lambda) \simeq 0\text{)}
\]
on a $R^{\flat} \subset S^{\flat} \times S^{\flat}$ formé des \struck{\ill{}} \add{couples} $(x^{\flat}, y^{\flat})$, avec $x, y \in M^{*}$, $x \wedge y \in (\det)^{*}$ --- i.e. les couples de positions des fibres projectifs distinctes en tout pt de la fibre. Le $S$-torseur relatif, sur \struck{$\mathbb{P}(M)$} $S^{\flat}$ correspond au fibré des \uncertain{semi-droites} \struck{\ill{}} \add{d)} et \uncertain{à la} \uncertain{multiplication} par $\pm 1$. L'iso. \struck{\ill{}} d) est induit de $\det M^{*} \simeq D^{*} \wedge_{\Lambda^{*}} D'^{*}$ en formant : $\Lambda^{*}/\pm 1$.
\marginal{\ill{} en \uncertain{considérant} les \uncertain{couples} $D, D'$ de \uncertain{droites} de $M$ et \ill{} $M = D \oplus D'$, d'où $\det M \simeq D \otimes D'$, donc $\det M^{*} \simeq D^{*} \wedge_{\Lambda^{*}} D'^{*}$}
\note{la fin du paragraphe et la note marginale (oblique) se lisent mal ; « $S$-torseur » et « $\check{S}^{*}$ » tels qu'écrits.}

\page{235}
\note{p. 120' de l'auteur (le chiffre des unités est surchargé). Au bas du feuillet, sous la bande collée, reparaît le haut de la feuille « 119 bis » (page précédente), non transcrit à nouveau.}

\section*{2. Condition d'isomorphie pour une donnée $\pi$.}

\struck{On note \ill{}}
\[
\text{\struck{$(45) \quad G = \{g \in \mathcal{A} = \mathrm{Aut}(\Gamma) \mid g \text{ centralise } Z \text{ i.e. } g \in \mathcal{G},\ g(\pi) = \ldots\}$}}
\]
On s'intéresse aux automorphismes d'une structure \struck{\ill{}}
\[
(49) \qquad \underset{\text{« prégraphe multijaugé »}}{\mathcal{M}} = (\underset{\substack{\| \\ (S, \check{S}, R)}}{\Gamma},\ Z \subset \mathcal{A} = \mathrm{Aut}(\Gamma),\ \pi \subset \sqcap_Z \subset \mathrm{Bij}(S, \check{S})) ,
\]
\note{sous « prégraphe multijaugé », un mot biffé.}
ce sont donc les $g \in \mathcal{A}$ qui \emph{normalisent} $Z$, et qui satisfont
\[
(50) \qquad \underset{\text{tr. de str.}}{g(\pi)} \Bigl( \overset{\mathrm{def}}{=} \underbrace{\{g_{\check{S}} \circ \varphi \circ g_S^{-1} \mid \varphi \in \pi\}}_{} \Bigr) = \pi
\]
ce qui équivaut à la condition, \struck{\ill{}} \add{en termes d'un $\varphi_0 \in \pi$}
\[
(51) \qquad g(\varphi_0) \Bigl( \overset{\mathrm{def}}{=} g_{\check{S}}\, \varphi_0\, g_S^{-1} \Bigr) \in \pi .
\]
On a donc des homs de groupes
\[
(52) \qquad
\begin{array}{c}
\mathrm{Aut}(\mathcal{M}) \longrightarrow \mathcal{A}/\!/_{Z}(\pi) \longrightarrow \mathrm{Aut}(Z) , \\
\cap \\
\mathcal{A}
\end{array}
\]
\note{« $\mathcal{A}/\!/_{Z}(\pi)$ » rend tel quel le signe écrit ; lecture incertaine.}
et \uncertain{par} ce qui \uncertain{intéresse} \add{aux automorphismes} de $\mathcal{M}$ dont l'image dans $\mathrm{Aut}(Z)$ sont égaux à $1$, i.e. dont l'image des $Z$ qui \emph{centralisent} $Z$ i.e. qui opèrent sur $\pi$ par translation, \struck{\ill{}} donc --- s'identifient --- groupe formé des $g$ :
\[
(53) \qquad G = \Bigl\{ g \in \mathcal{G} \overset{\mathrm{def}}{=} \mathrm{Centr}_{\mathcal{A}}(Z) \ \Bigm|\ 
\begin{array}{l}
g(\pi) = \pi \text{ i.e.} \\
g_{\check{S}}\, \varphi_0\, g_S^{-1} = \varphi_0
\end{array}
\Bigr\} .
\]
\note{la condition de droite est telle qu'écrite ($= \varphi_0$, et non $\in \pi$). La phrase qui précède est rapide et se lit mal.}

On trouve donc un diagramme canonique
\note{ces mots, au-dessus d'un double trait, sont la dernière ligne du feuillet ; la suite manque ici.}

\page{236}
\note{p. 121 de l'auteur ; suite directe de la p. 120' (« On trouve donc un diagramme canonique »).}
(54)

\begin{tikzcd}
1 \arrow[r] & G_0 \arrow[r] & G \arrow[r, "\delta"] & Z \\
 & & Z \arrow[u, hook] \arrow[ur, "\lambda \mapsto \lambda^{2}"'] &
\end{tikzcd}

où $\delta$ (jouant le rôle d'une « \uncertain{déterminant} ») exprime l'opération de $G$ sur $\pi$ par translations, et où on a posé
\[
(55) \qquad
\begin{aligned}
G_0 = \mathrm{Ker}\,\delta &= \{g \in \mathcal{A} \mid g(\varphi) = \varphi \ \ \forall \varphi \in \pi\} \\
&= \Bigl\{ g \in \mathcal{A} \ \Bigm|\ 
\begin{array}{l}
g \in \mathcal{G} \text{ i.e. } g \text{ centralise } Z \\
g(\varphi_0) = \varphi_0
\end{array}
\Bigr\} .
\end{aligned}
\]
\note{dans la première ligne, « $\mathcal{G}$ » biffé et remplacé par $\mathcal{A}$ ; au-dessus de « $g \in \mathcal{G}$ », un mot ajouté illisible.}
\marginal{NB Si $R$ séparé en $S^{\vee}$, … (55') $G_0 \xrightarrow{\ \sim\ }$ groupe des automorphismes du graphe $\gamma_{\varphi_0} = (S, A_{\varphi_0})$ qui commutent à $Z_S$, $\lambda \mapsto \lambda_S$.}
Le fait que $Z \subset G$, et la commutativité exprimée dans (54)
\[
(56) \qquad \delta(\lambda) = \lambda^{2} \qquad \forall \lambda \in Z
\]
provient de la relation
\[
\lambda(\varphi) = \underbrace{\lambda_{\check{S}} \varphi}_{\lambda \varphi} \underbrace{\lambda_S^{-1}}_{} = \lambda_{\check{S}}^{2} \varphi = T_{\lambda^{2}}(\varphi) .
\]
La définition de $\delta(g)$ pour $g \in G$ s'explicite par
\[
(57) \qquad g(\varphi) = T_{\delta(g)} \varphi \ \bigl(= \delta(g)_{\check{S}} \circ \varphi = \varphi \circ \delta(g)_S\bigr) \qquad \forall \varphi \in \pi,\ g \in G
\]
qui caractérise $\delta(g)$, pour \uncertain{un} \uncertain{seul} $\varphi = \varphi_0$ fixé déjà.

On a les
\[
(58) \qquad
\begin{array}{l}
\text{conditions} \\
\text{équivalentes} \\
\text{sur } \mathcal{M} = (\Gamma, Z, \pi)
\end{array}
\quad
\left\{
\begin{array}{l}
\text{a) L'hom } \delta \text{ est surjectif} \\
\text{b) \struck{\ill{}} $G$ opère transitivement sur } \pi \\
\text{c) } \forall \lambda \in Z, \text{ \struck{\ill{}} la structure de graphe définie par} \\
\qquad A_{\lambda} \subset \mathfrak{P}_{\{1,2\}}(S), \\
\qquad A_{\lambda} = \{\{s, \lambda t\} \mid s, t \in S,\ \text{t.q. } \{s,t\} \in A_0\} \\
\qquad \text{est isomorphe à } \gamma(S, A_0), \\
\qquad \text{par un } f \in \mathfrak{S}_S \text{ qui centralise } Z_S \\
\text{d) } \forall \lambda \in Z,\ \exists f_{\lambda} \in \mathfrak{S}_S, \text{ centralisant } Z_S, \text{ telle que} \\
\qquad \forall s, t \in S, \ \ldots \\
\qquad \{s,t\} \in A_0 \iff \{f_{\lambda}^{-1} s, f_{\lambda}^{-1} \lambda t\} \in A_0
\end{array}
\right.
\]
\marginal{ici on suppose choisi un $\varphi_0 \in \pi$, d'où une structure de graphe --- et on suppose $R$ séparé en $S^{\vee}$ \quad $\to$ d)}
\note{« $\forall \lambda \in Z$ » de d) est ajouté au-dessus d'un mot biffé.}

On dit que la condition d'homogénéité est satisfaite
\note{la phrase se poursuit p. 122 de l'auteur.}

\page{237}
\note{verso, sans numéro d'auteur : la phrase de la p. 121 s'y achève, puis le feuillet reste blanc.}
si l'hom $\delta$ est surjectif, i.e. les conditions \uncertain{(58)} satisfaites.

\struck{On}

\page{238}
\note{p. 122 de l'auteur.}

\section*{3. Conditions d'homogénéité.}

\struck{Nous nous intéressons maintenant au cas où} \struck{Revenons au cas de l'homogénéité} \struck{$G_0$ a la propriété d'homogénéité}
\marginal{Cas d'un}
\[
\mathcal{M} = (\Gamma,\ Z,\ \pi)
\]
\struck{\ill{}} \uncertain{Supposons} des conditions équivalentes entre les propriétés
\[
(59) \qquad
\left\{
\begin{array}{l}
\text{a) } G_0 \text{ est transitif sur } \overrightarrow{R} \\
\text{b) } G_0 \text{ est transitif sur } R_{\varphi_0} \text{ i.e. sur l'ens } \overrightarrow{A_{\varphi_0}} \\
\qquad \text{des arêtes orientées du graphe } (S, A_{\varphi_0}), \text{ pour} \\
\qquad \text{\emph{une} } \varphi_0 \in \pi \\
\text{b') comme b), pour tout } \varphi \in \pi
\end{array}
\right.
\]
\note{à gauche de l'accolade, un signe « $\lambda$ » ou « $\Lambda$ » isolé ; la flèche au-dessus de $R$ en a) est d'une encre plus tardive.}
\marginal{NB les \uncertain{conditions} a), b), b') \uncertain{signifient} aussi \ill{} --- $G_0$ transitif \ill{} \ill{} \ill{} graphes \ill{} $Z$. (Mais \ill{} \ill{} $G_0$ \ill{} \ill{} \ill{})}
\note{note marginale en oblique, très serrée, en partie biffée ; lecture fragmentaire.}
--- l'équivalence de a) et b) résulte du \struck{\ill{}} fait que
\[
R \simeq R_{\varphi_0} \quad (G_0\text{-isom}) \text{ via } \mathrm{id} \times \varphi_0 .
\]
Nous dirons que la situation \struck{du multigraphe} $\mathcal{M}$ \add{\uncertain{satisfait}} \struck{\ill{} $R \neq \emptyset$}, \struck{cette condition} \add{\ill{}} est \emph{homogène} s'il satisfait (59) et si $G_0$ est transitif sur $S$ et sur $S^{\vee}$ \struck{est \ill{}}.
\note{la fin de la phrase est surchargée de plusieurs corrections superposées.}

\struck{\ill{}} alors un $\lambda \in Z$ \struck{\ill{}} \uncertain{ayant} un pt fixe dans $S$ (resp. dans $S^{\vee}$), il opère trivialement dans $S$ (resp. $S^{\vee}$), donc il est trivial si $\Gamma$ est séparé en $S^{\vee}$ (resp. en $S$).

Donc
\[
(60) \qquad
\begin{array}{l}
\text{si } \Gamma \text{ séparé en } \check{S} \text{ (resp } S^{*}\text{), alors sous la condition} \\
\text{d'homogénéité (55), } Z \text{ opère librement sur } S \\
\text{(resp. } S^{\vee}\text{).}
\end{array}
\]
\note{« (55) » tel qu'écrit, ici et en (61) ; on attend la condition d'homogénéité (59).}
On a aussi
\[
(61) \qquad
\begin{array}{l}
\text{(55) implique (41), et ses formes équivalentes} \\
\text{(41') (41'') (44) (48) --- \uncertain{données dans les}} \\
\text{\uncertain{interprétations} (48.1) à (48.8).}
\end{array}
\]
\struck{Ce qui s'écrit encore \ill{} \ill{} \ill{} 59 a)) \ill{} \ill{}, i.e. $\pi$, \ill{} termes \ill{} et \ill{} $G_0, Z \subset \mathcal{A}$ deux sous-groupes de $\mathcal{A} = \mathrm{Aut}\,\Gamma$ commutant l'un à l'autre. \ill{}, \ill{} la forme 41 : \ill{} $(s_0, s_0^{\vee}) \in R$ et $(s_0, \lambda \check{s}_0) \in R$, $\lambda \in Z \implies \lambda = 1$ ; en effet il existe $g \in G_0$ tel que $g(s_0, \check{s}_0) = (s_0, \lambda \check{s}_0)$,}
\note{le dernier paragraphe est encadré et biffé de quatre traits obliques ; on le donne tel qu'il se lit.}

\page{239}
\note{verso sans numéro d'auteur (en tête, une abréviation : « \uncertain{Fin.} » ?) ; suite de la p. 122.}
\struck{$g s_0 = s_0$, $g \check{s}_0 = \lambda \check{s}_0$ i.e. $s_0 = g^{-1} \lambda s_0^{\vee} = \lambda g^{-1} s_0^{\vee}$}
en effet, comme $G_0$ opère transitivement sur \add{\ill{}} $A_0$ et $A_{\lambda}$ (il suffit \struck{\ill{}} \add{\ill{}} est supposé que l'un ait cette transitivité. D'ailleurs, au lieu de la transitivité sur $\overrightarrow{A_0}$ \struck{\ill{}} i.e. sur $R$), si on \struck{\ill{}} a $A_0 \cap A_{\lambda} \neq \emptyset$, on en conclut $A_0 = A_{\lambda}$ d'où $\lambda = \mathrm{id}$, cqfd.
\note{« $\neq \emptyset$ » : la barre du signe se lit mal ; la page porte peut-être « $= \emptyset$ ».}

Nous voulons maintenant expliciter la \struck{\ill{}} \add{\ill{}} \struck{donnée} d'un multigraphe \uncertain{homogène} (i.e. satisfaisant (59)), en termes de \uncertain{données} de théorie des groupes --- plus particulièrement dans le cas où on suppose en plus satisfait la condition d'isomorphie (58).

Pour ceci, \struck{dans le} \struck{\ill{}} je vais prendre les choses par un biais légèrement différent en partant d'un cas « épinglé » \struck{graphe} de \uncertain{départ} :
\[
(62) \qquad
\left\{
\begin{array}{l}
\text{a) graphe } \gamma_0 = (S, A_0), \text{ muni de \struck{\ill{}}} \\
\text{b) groupe \ill{} d'automorphismes de } (S, A_0), \text{ transitif sur } \overrightarrow{A_0} \\
\qquad \text{[et sur } S, \text{ i.e. } \forall s \in S,\ \mathrm{Et}_{A_0}(s) \neq \emptyset] \\
\text{c) et d'une « origine » } (s_0, t_0), \text{ \uncertain{avec} donc } \{s_0, t_0\} \in A_0 \\
\qquad \text{et enfin} \\
\text{d) d'un sous-groupe } Z \subset \mathfrak{S}_S \text{ commutatif,} \\
\qquad \text{formé d'homothéties pour le graphe } \gamma_0 \ \text{\ill{}} \\
\qquad \text{commutant \struck{\ill{}} : l'action de \ill{}}
\end{array}
\right.
\]
\note{le groupe de b) et de d) est noté d'un signe surchargé, corrigé en « $H$ » semble-t-il.}
\uncertain{Nous} nous préoccupons pas \uncertain{de} \uncertain{la} \uncertain{maximalité} de \ill{} \add{$H$}, \ill{} --- \uncertain{pour} le groupe $G_0$ formé de \emph{tous} les automorphismes \struck{qui} de la structure qui commutent à $Z$ --- \uncertain{et} \uncertain{il} \uncertain{suffit} \uncertain{j'aurai} \uncertain{aussi} \uncertain{à} \uncertain{travailler} avec le \uncertain{donnée} d'un sous-groupe plus petit, $G_0^{+}$. \struck{Autant dans \ill{} partir \ill{} \ill{} \ill{} $H$, sous-groupe, je vais le noter $H$, pour \ill{} \ill{} la donnée correspondante. \ill{} \ill{}}
\note{les dernières lignes sont rapides et se lisent mal ; le passage final est biffé de quatre traits.}

\page{240}
\note{p. 123 de l'auteur.}
\struck{la sous-structure de l'épinglage \ill{} dans un groupe $G$}
\add{\uncertain{vendredi}}
\note{ce mot, souligné, ajouté en tête près du numéro de page, est peut-être une date ; lecture très incertaine.}
Je considère \add{aussi} $H$ comme une structure supplémentaire, en plus de \struck{$Z$} groupe d'homothéties $Z$, et je définis $G$ en considérant
\[
(63) \qquad G = \text{\struck{groupe des}}\ \Bigl\{ f \in \mathfrak{S}_S \ \Bigm|\ 
\begin{array}{l}
f \text{ normalise } G, \text{ centralise } Z, \text{ et } \exists \lambda \in Z \text{ tel} \\
\text{que } f(A_0) = A_{\lambda} \text{ i.e.} \\
\forall s, t \in S, \ \{s,t\} \in A_0 \iff \{f^{-1}s, f^{-1}\lambda t\} \in A_0
\end{array}
\Bigr\}
\]
\note{« $f$ normalise $G$ » tel qu'écrit ; on attend $H$ (ou $G_0$).}
de sorte qu'on a encore un diagramme

(64)

\begin{tikzcd}
 & H \arrow[d, hook] \arrow[dr, dashed, "\text{distingué}"] & & \\
1 \arrow[r] & G_0 \arrow[r] & G \arrow[r, "\delta"] & Z \\
 & & Z \arrow[u, hook] \arrow[ur, "\lambda \mapsto \lambda^{2}"'] &
\end{tikzcd}

\note{sous $G_0$, « $= \mathrm{Ker}\,\delta$ ».}
\marginal{Pour écrire ici $Z$, et non un quotient de $Z$, on admet que $Z$ opère \emph{fidèlement} sur $\pi$, ce qui est automatique si $H$ \uncertain{commute} à $Z$ \uncertain{et} \uncertain{si} $\gamma_0$ est \ill{}, par \ill{}.}
où l'inclusion $H \to G_0 \to G$ fait de $H$ un sous-groupe \emph{distingué} de $G$. Pour l'instant on ne suppose pas $\delta$ surjectif. On fera \uncertain{abstraction} que dans cette présentation, même si $\delta$ surjectif, les groupes $G_0$, $G$ (\uncertain{intervenant} \add{ici}) \struck{\ill{}} \uncertain{peuvent} être plus petits que ceux des n°s précédents, qui eux ne \uncertain{dépendent} pas \uncertain{de} \uncertain{la} donnée préliminaire d'un $H$. En fait, ceux \struck{\ill{}} \add{\uncertain{envisagés} ici sont} les \emph{normalisateurs} de $H$ dans les groupes $G_0^{!}$, $G$ envisagés aux n°s précédents. L'inconvénient \struck{\ill{}} \uncertain{est} \uncertain{conséquence}, bien sûr, on \uncertain{prendra} $H = G_0$.

Notre objet principal est de décrire de façon \uncertain{aussi} précise que possible le diagramme (64),
\note{la phrase se poursuit au-delà du lot.}

\end{document}
